% Options for packages loaded elsewhere
\PassOptionsToPackage{unicode}{hyperref}
\PassOptionsToPackage{hyphens}{url}
\documentclass[
]{article}
\usepackage{xcolor}
\usepackage[margin=1in]{geometry}
\usepackage{amsmath,amssymb}
\setcounter{secnumdepth}{-\maxdimen} % remove section numbering
\usepackage{iftex}
\ifPDFTeX
  \usepackage[T1]{fontenc}
  \usepackage[utf8]{inputenc}
  \usepackage{textcomp} % provide euro and other symbols
\else % if luatex or xetex
  \usepackage{unicode-math} % this also loads fontspec
  \defaultfontfeatures{Scale=MatchLowercase}
  \defaultfontfeatures[\rmfamily]{Ligatures=TeX,Scale=1}
\fi
\usepackage{lmodern}
\ifPDFTeX\else
  % xetex/luatex font selection
\fi
% Use upquote if available, for straight quotes in verbatim environments
\IfFileExists{upquote.sty}{\usepackage{upquote}}{}
\IfFileExists{microtype.sty}{% use microtype if available
  \usepackage[]{microtype}
  \UseMicrotypeSet[protrusion]{basicmath} % disable protrusion for tt fonts
}{}
\makeatletter
\@ifundefined{KOMAClassName}{% if non-KOMA class
  \IfFileExists{parskip.sty}{%
    \usepackage{parskip}
  }{% else
    \setlength{\parindent}{0pt}
    \setlength{\parskip}{6pt plus 2pt minus 1pt}}
}{% if KOMA class
  \KOMAoptions{parskip=half}}
\makeatother
\usepackage{color}
\usepackage{fancyvrb}
\newcommand{\VerbBar}{|}
\newcommand{\VERB}{\Verb[commandchars=\\\{\}]}
\DefineVerbatimEnvironment{Highlighting}{Verbatim}{commandchars=\\\{\}}
% Add ',fontsize=\small' for more characters per line
\usepackage{framed}
\definecolor{shadecolor}{RGB}{248,248,248}
\newenvironment{Shaded}{\begin{snugshade}}{\end{snugshade}}
\newcommand{\AlertTok}[1]{\textcolor[rgb]{0.94,0.16,0.16}{#1}}
\newcommand{\AnnotationTok}[1]{\textcolor[rgb]{0.56,0.35,0.01}{\textbf{\textit{#1}}}}
\newcommand{\AttributeTok}[1]{\textcolor[rgb]{0.13,0.29,0.53}{#1}}
\newcommand{\BaseNTok}[1]{\textcolor[rgb]{0.00,0.00,0.81}{#1}}
\newcommand{\BuiltInTok}[1]{#1}
\newcommand{\CharTok}[1]{\textcolor[rgb]{0.31,0.60,0.02}{#1}}
\newcommand{\CommentTok}[1]{\textcolor[rgb]{0.56,0.35,0.01}{\textit{#1}}}
\newcommand{\CommentVarTok}[1]{\textcolor[rgb]{0.56,0.35,0.01}{\textbf{\textit{#1}}}}
\newcommand{\ConstantTok}[1]{\textcolor[rgb]{0.56,0.35,0.01}{#1}}
\newcommand{\ControlFlowTok}[1]{\textcolor[rgb]{0.13,0.29,0.53}{\textbf{#1}}}
\newcommand{\DataTypeTok}[1]{\textcolor[rgb]{0.13,0.29,0.53}{#1}}
\newcommand{\DecValTok}[1]{\textcolor[rgb]{0.00,0.00,0.81}{#1}}
\newcommand{\DocumentationTok}[1]{\textcolor[rgb]{0.56,0.35,0.01}{\textbf{\textit{#1}}}}
\newcommand{\ErrorTok}[1]{\textcolor[rgb]{0.64,0.00,0.00}{\textbf{#1}}}
\newcommand{\ExtensionTok}[1]{#1}
\newcommand{\FloatTok}[1]{\textcolor[rgb]{0.00,0.00,0.81}{#1}}
\newcommand{\FunctionTok}[1]{\textcolor[rgb]{0.13,0.29,0.53}{\textbf{#1}}}
\newcommand{\ImportTok}[1]{#1}
\newcommand{\InformationTok}[1]{\textcolor[rgb]{0.56,0.35,0.01}{\textbf{\textit{#1}}}}
\newcommand{\KeywordTok}[1]{\textcolor[rgb]{0.13,0.29,0.53}{\textbf{#1}}}
\newcommand{\NormalTok}[1]{#1}
\newcommand{\OperatorTok}[1]{\textcolor[rgb]{0.81,0.36,0.00}{\textbf{#1}}}
\newcommand{\OtherTok}[1]{\textcolor[rgb]{0.56,0.35,0.01}{#1}}
\newcommand{\PreprocessorTok}[1]{\textcolor[rgb]{0.56,0.35,0.01}{\textit{#1}}}
\newcommand{\RegionMarkerTok}[1]{#1}
\newcommand{\SpecialCharTok}[1]{\textcolor[rgb]{0.81,0.36,0.00}{\textbf{#1}}}
\newcommand{\SpecialStringTok}[1]{\textcolor[rgb]{0.31,0.60,0.02}{#1}}
\newcommand{\StringTok}[1]{\textcolor[rgb]{0.31,0.60,0.02}{#1}}
\newcommand{\VariableTok}[1]{\textcolor[rgb]{0.00,0.00,0.00}{#1}}
\newcommand{\VerbatimStringTok}[1]{\textcolor[rgb]{0.31,0.60,0.02}{#1}}
\newcommand{\WarningTok}[1]{\textcolor[rgb]{0.56,0.35,0.01}{\textbf{\textit{#1}}}}
\usepackage{graphicx}
\makeatletter
\newsavebox\pandoc@box
\newcommand*\pandocbounded[1]{% scales image to fit in text height/width
  \sbox\pandoc@box{#1}%
  \Gscale@div\@tempa{\textheight}{\dimexpr\ht\pandoc@box+\dp\pandoc@box\relax}%
  \Gscale@div\@tempb{\linewidth}{\wd\pandoc@box}%
  \ifdim\@tempb\p@<\@tempa\p@\let\@tempa\@tempb\fi% select the smaller of both
  \ifdim\@tempa\p@<\p@\scalebox{\@tempa}{\usebox\pandoc@box}%
  \else\usebox{\pandoc@box}%
  \fi%
}
% Set default figure placement to htbp
\def\fps@figure{htbp}
\makeatother
\setlength{\emergencystretch}{3em} % prevent overfull lines
\providecommand{\tightlist}{%
  \setlength{\itemsep}{0pt}\setlength{\parskip}{0pt}}
\usepackage{bookmark}
\IfFileExists{xurl.sty}{\usepackage{xurl}}{} % add URL line breaks if available
\urlstyle{same}
\hypersetup{
  pdftitle={GLMMs and GEE},
  hidelinks,
  pdfcreator={LaTeX via pandoc}}

\title{GLMMs and GEE}
\author{}
\date{\vspace{-2.5em}}

\begin{document}
\maketitle

\textbf{Inference lab} using the five-step template: Hypothesis →
Visualise → Assumptions → Conduct → Conclude.

\subsection{Learning objectives}\label{learning-objectives}

\begin{itemize}
\tightlist
\item
  Fit a binary GLMM for clustered data with \texttt{lme4::glmer} and
  with \texttt{glmmTMB}.
\item
  Fit a GEE with \texttt{geepack::geeglm} and compare coefficient
  interpretations.
\item
  Say why GLMM and GEE estimate different things (subject-specific vs
  population-average).
\end{itemize}

\subsection{Prerequisites}\label{prerequisites}

Linear mixed models and logistic regression.

\subsection{Background}\label{background}

For binary outcomes in clustered data, the two standard approaches are
generalised linear mixed models (GLMM) and generalised estimating
equations (GEE). A GLMM is a random-effects model on the link scale; its
fixed-effect coefficients are conditional on the random effect and
describe subject-specific effects. A GEE fits a marginal model using a
working correlation structure; its coefficients describe
population-averaged effects. For logistic models, the two will not in
general agree even when both are correctly specified.

\texttt{lme4::glmer} is the classic random-effects fitter;
\texttt{glmmTMB} handles larger and more complex random structures
including correlated random effects and zero-inflated distributions.
\texttt{geepack} fits GEEs with several working correlation choices; the
standard errors are robust to misspecification of the correlation
structure but the point estimate can shift.

The choice between GLMM and GEE is not a statistical technicality; it is
a choice of estimand. If you care about what happens inside a person
over time, you want a GLMM. If you care about a population- average
effect for policy or for a marginal treatment effect, you want a GEE (or
a marginalised GLMM).

\subsection{Setup}\label{setup}

\begin{Shaded}
\begin{Highlighting}[]
\FunctionTok{library}\NormalTok{(tidyverse)}
\FunctionTok{library}\NormalTok{(lme4)}
\FunctionTok{library}\NormalTok{(glmmTMB)}
\FunctionTok{library}\NormalTok{(geepack)}
\FunctionTok{set.seed}\NormalTok{(}\DecValTok{42}\NormalTok{)}
\FunctionTok{theme\_set}\NormalTok{(}\FunctionTok{theme\_minimal}\NormalTok{(}\AttributeTok{base\_size =} \DecValTok{12}\NormalTok{))}
\end{Highlighting}
\end{Shaded}

\subsection{1. Hypothesis}\label{hypothesis}

Treatment reduces the probability of a repeated binary event in a
clustered study (clusters = clinics). The true log-odds fixed effect is
−0.5.

\subsection{2. Visualise}\label{visualise}

\begin{Shaded}
\begin{Highlighting}[]
\NormalTok{nj }\OtherTok{\textless{}{-}} \DecValTok{20}
\NormalTok{dat }\OtherTok{\textless{}{-}} \FunctionTok{tibble}\NormalTok{(}
  \AttributeTok{clinic =} \FunctionTok{factor}\NormalTok{(}\FunctionTok{rep}\NormalTok{(}\FunctionTok{seq\_len}\NormalTok{(J), }\AttributeTok{each =}\NormalTok{ nj)),}
  \AttributeTok{trt    =} \FunctionTok{rbinom}\NormalTok{(J }\SpecialCharTok{*}\NormalTok{ nj, }\DecValTok{1}\NormalTok{, }\FloatTok{0.5}\NormalTok{),}
  \AttributeTok{u      =} \FunctionTok{rep}\NormalTok{(}\FunctionTok{rnorm}\NormalTok{(J, }\DecValTok{0}\NormalTok{, }\FloatTok{0.8}\NormalTok{), }\AttributeTok{each =}\NormalTok{ nj)}
\NormalTok{) }\SpecialCharTok{|\textgreater{}}
  \FunctionTok{mutate}\NormalTok{(}\AttributeTok{lp =} \SpecialCharTok{{-}}\FloatTok{0.5} \SpecialCharTok{*}\NormalTok{ trt }\SpecialCharTok{+}\NormalTok{ u,}
         \AttributeTok{y  =} \FunctionTok{rbinom}\NormalTok{(}\FunctionTok{n}\NormalTok{(), }\DecValTok{1}\NormalTok{, }\FunctionTok{plogis}\NormalTok{(lp)))}

\NormalTok{dat }\SpecialCharTok{|\textgreater{}}
  \FunctionTok{group\_by}\NormalTok{(clinic, trt) }\SpecialCharTok{|\textgreater{}}
  \FunctionTok{summarise}\NormalTok{(}\AttributeTok{p =} \FunctionTok{mean}\NormalTok{(y), }\AttributeTok{.groups =} \StringTok{"drop"}\NormalTok{) }\SpecialCharTok{|\textgreater{}}
  \FunctionTok{ggplot}\NormalTok{(}\FunctionTok{aes}\NormalTok{(}\FunctionTok{factor}\NormalTok{(trt), p)) }\SpecialCharTok{+}
  \FunctionTok{geom\_boxplot}\NormalTok{(}\AttributeTok{alpha =} \FloatTok{0.6}\NormalTok{) }\SpecialCharTok{+}
  \FunctionTok{labs}\NormalTok{(}\AttributeTok{x =} \StringTok{"Treatment"}\NormalTok{, }\AttributeTok{y =} \StringTok{"Event probability per clinic"}\NormalTok{)}
\end{Highlighting}
\end{Shaded}

\subsection{3. Assumptions}\label{assumptions}

GLMM: normal random intercepts for clinic, correct link. GEE: correct
mean structure; standard errors robust to working correlation.

\subsection{4. Conduct}\label{conduct}

\begin{Shaded}
\begin{Highlighting}[]
\NormalTok{                   data }\OtherTok{=}\NormalTok{ dat, family }\OtherTok{=}\NormalTok{ binomial}\ErrorTok{)}
\NormalTok{fit\_tmb   }\OtherTok{\textless{}{-}} \FunctionTok{glmmTMB}\NormalTok{(y }\SpecialCharTok{\textasciitilde{}}\NormalTok{ trt }\SpecialCharTok{+}\NormalTok{ (}\DecValTok{1} \SpecialCharTok{|}\NormalTok{ clinic),}
                     \AttributeTok{data =}\NormalTok{ dat, }\AttributeTok{family =}\NormalTok{ binomial)}
\NormalTok{fit\_gee   }\OtherTok{\textless{}{-}} \FunctionTok{geeglm}\NormalTok{(y }\SpecialCharTok{\textasciitilde{}}\NormalTok{ trt, }\AttributeTok{id =}\NormalTok{ clinic, }\AttributeTok{data =}\NormalTok{ dat,}
                    \AttributeTok{family =}\NormalTok{ binomial, }\AttributeTok{corstr =} \StringTok{"exchangeable"}\NormalTok{)}

\FunctionTok{summary}\NormalTok{(fit\_glmer)}\SpecialCharTok{$}\NormalTok{coefficients}
\FunctionTok{summary}\NormalTok{(fit\_tmb)}\SpecialCharTok{$}\NormalTok{coefficients}\SpecialCharTok{$}\NormalTok{cond}
\FunctionTok{summary}\NormalTok{(fit\_gee)}\SpecialCharTok{$}\NormalTok{coefficients}
\end{Highlighting}
\end{Shaded}

\begin{Shaded}
\begin{Highlighting}[]
\NormalTok{  glmer   }\OtherTok{=}\NormalTok{ broom.mixed}\SpecialCharTok{::}\FunctionTok{tidy}\NormalTok{(fit\_glmer, }\AttributeTok{effects =} \StringTok{"fixed"}\NormalTok{),}
\NormalTok{  glmmTMB }\OtherTok{=}\NormalTok{ broom.mixed}\SpecialCharTok{::}\FunctionTok{tidy}\NormalTok{(fit\_tmb,   }\AttributeTok{effects =} \StringTok{"fixed"}\NormalTok{),}
\NormalTok{  gee     }\OtherTok{=}\NormalTok{ broom}\SpecialCharTok{::}\FunctionTok{tidy}\NormalTok{(fit\_gee),}
\NormalTok{  .id }\OtherTok{=} \StringTok{"model"}
\ErrorTok{)} \SpecialCharTok{|\textgreater{}}
  \FunctionTok{filter}\NormalTok{(term }\SpecialCharTok{==} \StringTok{"trt"}\NormalTok{) }\SpecialCharTok{|\textgreater{}}
  \FunctionTok{select}\NormalTok{(model, estimate, std.error)}
\end{Highlighting}
\end{Shaded}

\subsection{5. Concluding statement}\label{concluding-statement}

\begin{quote}
The GLMM gave a subject-specific treatment log-odds of about
\texttt{round(fixef(fit\_glmer){[}"trt"{]},\ 2)}; the GEE gave a
population- averaged log-odds of about
\texttt{round(coef(fit\_gee){[}"trt"{]},\ 2)}. The GEE estimate is
typically attenuated toward zero relative to the GLMM; both are correct
answers to different questions.
\end{quote}

\subsection{Common pitfalls}\label{common-pitfalls}

\begin{itemize}
\tightlist
\item
  Comparing GLMM and GEE coefficients as if they estimate the same
  thing.
\item
  Using a working independence structure in GEE when clusters are large
  and treatment is heterogeneous within them.
\item
  Ignoring overdispersion in a count GLMM by assuming Poisson when
  negative binomial fits better.
\item
  Fitting a GLMM with too few clusters (\textless{} 10) and trusting the
  variance component.
\end{itemize}

\subsection{Further reading}\label{further-reading}

\begin{itemize}
\tightlist
\item
  Liang KY, Zeger SL (1986), \emph{Longitudinal data analysis using
  generalized linear models}.
\item
  Zeger SL, Liang KY, Albert PS (1988), \emph{Models for longitudinal
  data: a generalized estimating equation approach}.
\item
  Brooks ME et al.~(2017), \emph{glmmTMB balances speed and
  flexibility}.
\end{itemize}

\subsection{Session info}\label{session-info}

\begin{Shaded}
\begin{Highlighting}[]

\end{Highlighting}
\end{Shaded}

\subsection{See also --- chapter index}\label{see-also-chapter-index}

\begin{itemize}
\tightlist
\item
  \href{../index.Rmd}{Methods in this chapter}
\item
  \href{../../../ch00-find-your-method.Rmd}{Find your method (Chapter
  0)}
\end{itemize}

\end{document}
